\begin{abstract}
\noindent Proofs show axiom instances, not axioms. We ask when a checker can learn from them which axioms and schemas a
community uses, soundly against adversarial provers. As in an earlier report, the checker is cautious: it accepts the
common instances of all templates in a fixed class that cover the data. We answer three questions.
(1)~Instances $\varphi(t)$ determine the instance schema $\varphi(z)$, a first-order pattern; for one variable, two with
different head symbols suffice. Passing to $\forall x\varphi$ is an $\omega$-rule step: truth-safe for all $\varphi$
exactly when the closed-term substructure is elementary (in $\mathbb N$, not in $\mathbb R$), and not derivable in
general.
(2)~Every ZF schema, in each formulation considered, is a higher-order pattern; from patterns to rigid second-order
templates, instances determine it iff their bodies' main symbols vary and every argument place is used. First-order
anti-unification works only for some formulations and encodings; for most others no finite union of first-order schemas
is sound. PA induction is not a pattern: anti-unification fails on it, determinate templates learn it.
(3)~For unlabelled mixtures we propose \DTRC{}: cluster the data by refuting merged templates, then check each cluster
cautiously. Under refutation separation it recovers the labels and needs no more data than a labelled learner; otherwise
it is sound relative to a residue of unrefuted merges, and dependence on refutation depth is unavoidable. We prove
separation for PA, check it on sampled ZF and ZFC data, and report experiments; in one, injected mistakes made \DTRC{}
accept a false sentence.
\end{abstract}
